\chapter{Introduction}\label{ch:introduction}

Industrial automation is evolving from traditional rule-based control towards
systems that incorporate artificial intelligence and, increasingly, agentic
behaviour~\cite{aslam2025industrial}. At the same time, the numbers of connected
IoT devices are growing, and each device spreads and widens the attack surface. Weak
authentication, unprotected update channels, and exposure of personal data are
recent recurring problems now~\cite{alotaibi2023iiot,meziane2023survey}. Standardization
bodies have responded with security baselines for these technologies. Two of
them are central to this thesis. \standardA{} defines baseline cybersecurity
provisions for consumer IoT devices, covering areas such as password policy,
secure communication, software updates, and security of personal
data~\cite{etsi303645}. \standardB{}, published by the ETSI committee on
Securing Artificial Intelligence, defines baseline security requirements for AI
models and systems across their life cycle, from secure design to secure end of
life~\cite{etsi304223}.

A company that relys on an AI-enabled IoT product falls under both documents.
Demonstrating compliance means working out how the two sets of
requirements related to each other. Some provisions say essentially the same thing, others
overlap only partially, or cover different aspects of one shared objective, and
many have no counterpart at all. Both standards express their requirements in
natural language, at different levels of abstraction and with
domain-specific vocabulary, this is the reason why mapping work is done manually today. It is
slow work, it is prone to error, and it requires highly skilled human in both IoT and AI
security~\cite{kosenkov2025systematic,djebbar2023comparative}. The burden grows
with every updates on regulation, as the EU Cyber Resilience Act and the AI Act
add their own requirement standards~\cite{canavese2026uncovering}.

Recent progress in natural language processing offers an idea to minimise this
effort. Embedding models built on transformers can successfully measure semantic similarity
between sentences even though the wording differs~\cite{devlin2019bert,
reimers2019sbert}, and popular large language models can be prompted to judge the
relationship between two requirements directly~\cite{rodriguez2024llm}. The thesis investigate on Whether these techniques are accurate enough for security standards, where precision and
traceability matter,which is still an open question. 

\section{Motivation}\label{sec:motivation}

Compliance mapping sits in an awkward middle position. It is very language-heavy
for classical rule-based tooling, yet too consequential to hand over to an
opaque model without evidence of its reliability. Published work on automated
compliance analysis reports promising results in neighbouring domains such as
privacy policies, GDPR data processing agreements, and cloud security
posture~\cite{rodriguez2024llm,cejas2023gdpr,barhaim2023cloud}, but the
specific task of mapping one security standard onto another has received little
systematic evaluation. By implementing a spectrum of methods, from keyword
matching to LLM classification, and scoring them against a common annotated
reference on a real standard pair, this project produces evidence about what
automation does well on this task, what it does badly, and what still needs a
human. The results are relevant both to researchers working on NLP for
requirements engineering and to practitioners who need to decide whether
AI-assisted compliance tools can be trusted in their processes.

\section{Aim}\label{sec:aim}

The project builds and evaluates an NLP-based pipeline that identifies semantic
relationships between the provisions of \standardA{} and \standardB{}. The aim
is not to develop new model architectures but to establish, empirically, how
well existing techniques perform on this task. Concretely, the project (i)
extracts the normative provisions of both standards into a structured dataset,
(ii) applies eleven automated mapping methods spanning rule-based, embedding-based,
and LLM-based approaches, (iii) evaluates them against a ground truth of annotated provision pairs that was
LLM-drafted, author-reviewed, and cross-verified by a second independent model in two of the three annotation rounds, and (iv) analyses the errors to
characterise the limitations of each method family. The outcome is a grounded
assessment of the feasibility of automated compliance mapping between IoT and
AI security standards, together with a reusable dataset and evaluation
framework.

\section{Research questions}\label{sec:rqs}

\begin{enumerate}[label=\textbf{RQ\arabic*.},leftmargin=3.2em]
  \item To what extent can NLP and AI techniques automatically identify
  semantic relationships (equivalence, overlap, subsumption, complementarity)
  between security requirements in \standardA{} and \standardB{}, and what are
  the fundamental challenges and limitations?
  \item How do different automated approaches compare on the compliance mapping
  task: lexical baselines (TF--IDF, security-keyword Jaccard),
  sentence-embedding models (Sentence-BERT, MPNet, BGE), contextual-embedding
  models (BERT, SecureBERT, CySecBERT), a cross-encoder trained for natural
  language inference, and hosted API methods (Gemini embeddings and Gemini
  prompt classification)?
  \item How does automated mapping compare to human judgement in terms of
  accuracy, coverage, and efficiency? The comparison is made against the
  annotated reference set of this project, which was LLM-drafted, cross-verified
  and author-reviewed rather than built by external domain experts; what that
  substitution costs the answer is set out in Sections~\ref{sec:gt}
  and~\ref{sec:threats}.
\end{enumerate}

\section{Contributions}\label{sec:contributions}

The main contributions of this degree project are:

\begin{itemize}
  \item The project provides an in-depth analysis of NLP approaches on automated compliance on ETSI EN 303 645 and EN 304 223 standards. A structured dataset of 141 normative provisions extracted from
  \standardA{} and \standardB{}, and a reference mapping of 200 pairs annotated
  with six relationship types, justifications, and confidence scores, drawn from
  1{,}027 pairs judged across three annotation rounds and balanced so that the
  trivial ``everything is related'' classifier does not score near the top of
  the scale.
  \item The project uses an open, reproducible pipeline implementing eleven mapping methods
  (TF--IDF, security-keyword Jaccard matching, Sentence-BERT, MPNet, BGE, BERT,
  SecureBERT, CySecBERT, an NLI cross-encoder, Gemini embeddings, and Gemini LLM
  classification) behind a common evaluation harness.
  \item The project presents a empirical comparison of these methods on pair detection and
  relationship classification, including confidence intervals, threshold
  sensitivity analysis, and qualitative error analysis, together with two
  controls that separate the methods from their setup: an ablation that
  re-scores every local method on normative text alone, and a classifier trained
  on the judged pairs held out of the evaluation.
  \item The project highlights the Evidence-based observations on why general-purpose embedding models
  overestimate similarity between security requirements, why cybersecurity-specific
  pretraining did not help in either independent attempt tested, and on the
  practical role automated mapping can play in compliance work today.
  \item The project helps the finding that equivalence and subsumption are close to absent between
  these two standards: 13 pairs among the 1{,}027 judged, and none among 300
  top-ranked pairs searched specifically to find more. This bounds what any typed
  mapping between a device-level and a lifecycle-level standard can report.
  \item The project helps to contribute the research community by supporting further research in automated compliance, cybersecurity, and AI governance.
\end{itemize}

\section{Outline}

Chapter~\ref{ch:background} introduces the two standards and reviews related
work on automated compliance analysis. Chapter~\ref{ch:methodology} describes
the research design: data extraction, the relationship taxonomy, ground-truth
construction and validation, the mapping methods, and the evaluation metrics.
Chapter~\ref{ch:prototype} presents the implemented pipeline.
Chapter~\ref{ch:results} reports and discusses the experimental results, and
Chapter~\ref{ch:conclusion} concludes and outlines future work.
